% manuale beamer http://www.ctan.org/tex-archive/macros/latex/contrib/beamer/doc/beameruserguide.pdf

% ==================///==================///==================///
% ==================/// LATEX'S STUFF
% ==================///==================///==================///
% !TEX root = main.tex
\documentclass[aspectratio={169}]{beamer}
\usepackage{amsfonts,amsmath,oldgerm}
\usetheme{_statale}
\usefonttheme[onlymath]{serif}

% COMANDI CUSTOM DEL TEMA:
% \testcolor{color}      — colorbox per visualizzare/verificare un colore
% \hrefcol{url}{testo}   — hyperlink colorato in cyan
% \course{...}           — campo corso di laurea (comando del tema _statale)
% \IDnumber{...}         — numero di matricola (campo del tema _statale)
% \footlinecolor{color}  — colore barra inferiore (switchabile per sezione)
%   colori disponibili: maincolor, stataledarkgreen, statalegreen,
%   statalelightgreen, statalered, stataleyellow, statalelilla
% \backmatter            — slide finale di ringraziamenti
% \backmatter[notitle]   — variante senza ripetizione del titolo
% Vedere chapters/special_slides.tex per gli ambienti chapter e sidepic.
\newcommand{\testcolor}[1]{\colorbox{#1}{\textcolor{#1}{test}}~\texttt{#1}}
\newcommand{\hrefcol}[2]{\textcolor{cyan}{\href{#1}{#2}}}
\titlebackground*{assets/background}

\usepackage{tikz}
\usetikzlibrary{arrows.meta, backgrounds, calc, chains, fit, graphs, graphdrawing, graphdrawing.layered, intersections, positioning, shadings, shapes, shapes.geometric, shapes.misc}

% STILI TIKZ GLOBALI (disponibili in tutti i frame):
% base        — stile base: testo centrato, inner sep 0.5em
% boxes       — rettangolo con bordo nero
% rboxes      — boxes con angoli arrotondati
% arrow_label — etichetta bianca centrata sulle frecce (usare come nodo su un path)
% Stili avanzati (entity, artifact, service, highlight*, lattice*, connector*):
% vedere 00_common_tikz.tex
\tikzset{%
	base/.style = {
			% minimum width=1cm,
			% minimum height=1cm,
			text centered,
			% font=\sffamily,
			inner sep=0.5em,
		},
	boxes/.style = {
			base,
			draw=black,
			rectangle,
		},
	rboxes/.style = {
			boxes,
			rounded corners
		},
	arrow_label/.style = {
			base,
			midway,
			% sloped,
			fill=white, anchor=center,
			inner sep=0.25em
		}
}

% for appendix
\usepackage{appendixnumberbeamer}

% ==================///==================///==================///
% ==================/// SPLASH PAGE
% ==================///==================///==================///

% !TEX root = ../main.tex
% Metadati presentazione — modifica questo file per ogni nuova presentazione.

\title{TODO: titolo tesi magistrale} % TODO
\course{Laurea Magistrale in TODO: corso di laurea} % TODO
\author{\href{mailto:TODO: email@studenti.unimi.it}{TODO: Nome Cognome}} % TODO
\IDnumber{TODO: matricola} % TODO
\date{
	Relatore: Prof. TODO: Nome Relatore\\ % TODO
	Correlatore: TODO: Titolo Nome Correlatore\\ % TODO
	TODO: data discussione} % TODO


% ==================///==================///==================///
% ==================/// START PRESENTATION
% ==================///==================///==================///

\begin{document}
% !TEX root = main.tex
% STILI TIKZ AVANZATI — disponibili in tutti i frame dopo % !TEX root = main.tex
% STILI TIKZ AVANZATI — disponibili in tutti i frame dopo % !TEX root = main.tex
% STILI TIKZ AVANZATI — disponibili in tutti i frame dopo \input{00_common_tikz.tex}
%
% ENTITÀ E STRUTTURE:
%   entity          — ellisse sottile (nodi concettuali)
%   artifact        — rettangolo arrotondato con sfondo grigio
%   service         — rettangolo smussato a 30° (chamfered rectangle)
%
% EVIDENZIAZIONE (colore oro/LightGoldenRod):
%   highlight       — bordo spesso oro con riempimento
%   highlight profile — riempimento LightGoldenRod!100
%   highlight double — doppio bordo oro
%   highlight text  — testo Goldenrod
%   highlight link  — link Goldenrod
%   highlight light — riempimento 50%
%   highlight very light — riempimento 20%
%
% CONNETTORI:
%   connector       — frecce Stealth piene
%   connector assurance — frecce Stealth tratteggiate
%   link            — frecce Latex
%   link comment    — linee tratteggiate
%   on link         — nodo bianco su un path
%
% RETICOLO (per diagrammi lattice/Hasse):
%   lattice layer label, lattice layer separator
%   lattice element bound, lattice element selected
%   lattice element not chosen, lattice element not chosen link
%
% FRECCE PERSONALIZZATE:
%   arrow           — shape=signal DarkBlue/White (stile "pipeline")
%   desc            — rettangolo VeryLightGrey con font piccolo
%
% OVERLAY CONDIZIONALE:
%   on slide/.code args={<#1>#2} — mostra un nodo solo su una slide specifica
%     es: \node[boxes, on slide=<2>{opacity=0}] {testo};
%
% MACRO:
%   \highlightprofilecolor — colore corrente highlight (ridefinibile)
\newcommand{\highlightprofilecolor}{LightGoldenRod!100}

\tikzset{
	entity/.style={
			draw,
			ellipse,
			very thin
		},
	artifact/.style={
			rectangle,
			draw,
			rounded corners,
			fill=VeryLightGrey!50
		},
	service/.style={
			chamfered rectangle,
			chamfered rectangle angle=30,
			draw,
			%fill=VeryLightGrey!120
		},
	every node/.style={
			font={\scriptsize},
			text centered,
		},
	on link/.style={
			fill=white,
		},
	connector/.style={
			>=stealth,
		},
	connector assurance/.style={
			>=Stealth,
			dashed
		},
	on slide/.code args={<#1>#2}{%
			\only<#1>{\pgfkeysalso{#2}}%
		},
	highlight base/.style={
			very thick
		},
	highlight/.style={
			highlight base,
			fill=LightGoldenRod,
		},
	highlight profile/.style={
			\highlightprofilecolor
		},
	highlight double/.style={
			double=LightGoldenRod
		},
	highlight text/.style={
			text=Goldenrod,
		},
	highlight link/.style={
			%thick,
			color=Goldenrod
		},
	highlight light/.style={
			fill=LightGoldenRod!50
		},
	highlight very light/.style={
			fill=LightGoldenRod!20
		},
	element/.style={
			inner sep=1.5pt,
			font=\footnotesize
		},
	connector/.style={
			>=Stealth,
		},
	lattice layer label/.style={
			text width=2.9cm,
			font={\footnotesize},
		},
	not chosen color/.style={
			black!50
		},
	lattice layer separator/.style={
			dashed,
		},
	service label/.style={
			fill=white,
			circle,
			draw,
			very thick,
			inner sep=2.5pt,
		},
	lattice element bound/.style={
			fill=black!15,
			double
		},
	lattice element selected/.style={
			fill=Goldenrod!15,
			double
		},
	lattice element not chosen/.style={
			not chosen color,
			densely dotted
		},
	lattice element not chosen link/.style={
			not chosen color,
			densely dotted
		},
	service element not chosen/.style={
			text=black!50,
			draw=black!50,
			dashed,
		},
	header/.style={
		},
	main box/.style={
			font={\small}
		},
	link/.style={
			>=Latex,
		},
	activity/.style={
			fill=white
		},
	% for cert history
	arrow/.style={
			draw,
			minimum height=1cm,
			inner sep=2pt,
			shape=signal,
			signal from=west,
			signal to=east,
			signal pointer angle=110,
			top color=DarkBlue,
			bottom color=DarkBlue,
			text=White,
			text width=.24\textwidth,
			font = {\small}
		},
	desc/.style={
			%text width=3.9cm,
			text width=.22\textwidth,
			fill=VeryLightGrey,
			%minimum height=1cm,
			%anchor=south
			% text depth=1cm,
			rectangle,
			font = {\footnotesize}
		},
	link comment/.style={
			dashed
		}
}
%
% ENTITÀ E STRUTTURE:
%   entity          — ellisse sottile (nodi concettuali)
%   artifact        — rettangolo arrotondato con sfondo grigio
%   service         — rettangolo smussato a 30° (chamfered rectangle)
%
% EVIDENZIAZIONE (colore oro/LightGoldenRod):
%   highlight       — bordo spesso oro con riempimento
%   highlight profile — riempimento LightGoldenRod!100
%   highlight double — doppio bordo oro
%   highlight text  — testo Goldenrod
%   highlight link  — link Goldenrod
%   highlight light — riempimento 50%
%   highlight very light — riempimento 20%
%
% CONNETTORI:
%   connector       — frecce Stealth piene
%   connector assurance — frecce Stealth tratteggiate
%   link            — frecce Latex
%   link comment    — linee tratteggiate
%   on link         — nodo bianco su un path
%
% RETICOLO (per diagrammi lattice/Hasse):
%   lattice layer label, lattice layer separator
%   lattice element bound, lattice element selected
%   lattice element not chosen, lattice element not chosen link
%
% FRECCE PERSONALIZZATE:
%   arrow           — shape=signal DarkBlue/White (stile "pipeline")
%   desc            — rettangolo VeryLightGrey con font piccolo
%
% OVERLAY CONDIZIONALE:
%   on slide/.code args={<#1>#2} — mostra un nodo solo su una slide specifica
%     es: \node[boxes, on slide=<2>{opacity=0}] {testo};
%
% MACRO:
%   \highlightprofilecolor — colore corrente highlight (ridefinibile)
\newcommand{\highlightprofilecolor}{LightGoldenRod!100}

\tikzset{
	entity/.style={
			draw,
			ellipse,
			very thin
		},
	artifact/.style={
			rectangle,
			draw,
			rounded corners,
			fill=VeryLightGrey!50
		},
	service/.style={
			chamfered rectangle,
			chamfered rectangle angle=30,
			draw,
			%fill=VeryLightGrey!120
		},
	every node/.style={
			font={\scriptsize},
			text centered,
		},
	on link/.style={
			fill=white,
		},
	connector/.style={
			>=stealth,
		},
	connector assurance/.style={
			>=Stealth,
			dashed
		},
	on slide/.code args={<#1>#2}{%
			\only<#1>{\pgfkeysalso{#2}}%
		},
	highlight base/.style={
			very thick
		},
	highlight/.style={
			highlight base,
			fill=LightGoldenRod,
		},
	highlight profile/.style={
			\highlightprofilecolor
		},
	highlight double/.style={
			double=LightGoldenRod
		},
	highlight text/.style={
			text=Goldenrod,
		},
	highlight link/.style={
			%thick,
			color=Goldenrod
		},
	highlight light/.style={
			fill=LightGoldenRod!50
		},
	highlight very light/.style={
			fill=LightGoldenRod!20
		},
	element/.style={
			inner sep=1.5pt,
			font=\footnotesize
		},
	connector/.style={
			>=Stealth,
		},
	lattice layer label/.style={
			text width=2.9cm,
			font={\footnotesize},
		},
	not chosen color/.style={
			black!50
		},
	lattice layer separator/.style={
			dashed,
		},
	service label/.style={
			fill=white,
			circle,
			draw,
			very thick,
			inner sep=2.5pt,
		},
	lattice element bound/.style={
			fill=black!15,
			double
		},
	lattice element selected/.style={
			fill=Goldenrod!15,
			double
		},
	lattice element not chosen/.style={
			not chosen color,
			densely dotted
		},
	lattice element not chosen link/.style={
			not chosen color,
			densely dotted
		},
	service element not chosen/.style={
			text=black!50,
			draw=black!50,
			dashed,
		},
	header/.style={
		},
	main box/.style={
			font={\small}
		},
	link/.style={
			>=Latex,
		},
	activity/.style={
			fill=white
		},
	% for cert history
	arrow/.style={
			draw,
			minimum height=1cm,
			inner sep=2pt,
			shape=signal,
			signal from=west,
			signal to=east,
			signal pointer angle=110,
			top color=DarkBlue,
			bottom color=DarkBlue,
			text=White,
			text width=.24\textwidth,
			font = {\small}
		},
	desc/.style={
			%text width=3.9cm,
			text width=.22\textwidth,
			fill=VeryLightGrey,
			%minimum height=1cm,
			%anchor=south
			% text depth=1cm,
			rectangle,
			font = {\footnotesize}
		},
	link comment/.style={
			dashed
		}
}
%
% ENTITÀ E STRUTTURE:
%   entity          — ellisse sottile (nodi concettuali)
%   artifact        — rettangolo arrotondato con sfondo grigio
%   service         — rettangolo smussato a 30° (chamfered rectangle)
%
% EVIDENZIAZIONE (colore oro/LightGoldenRod):
%   highlight       — bordo spesso oro con riempimento
%   highlight profile — riempimento LightGoldenRod!100
%   highlight double — doppio bordo oro
%   highlight text  — testo Goldenrod
%   highlight link  — link Goldenrod
%   highlight light — riempimento 50%
%   highlight very light — riempimento 20%
%
% CONNETTORI:
%   connector       — frecce Stealth piene
%   connector assurance — frecce Stealth tratteggiate
%   link            — frecce Latex
%   link comment    — linee tratteggiate
%   on link         — nodo bianco su un path
%
% RETICOLO (per diagrammi lattice/Hasse):
%   lattice layer label, lattice layer separator
%   lattice element bound, lattice element selected
%   lattice element not chosen, lattice element not chosen link
%
% FRECCE PERSONALIZZATE:
%   arrow           — shape=signal DarkBlue/White (stile "pipeline")
%   desc            — rettangolo VeryLightGrey con font piccolo
%
% OVERLAY CONDIZIONALE:
%   on slide/.code args={<#1>#2} — mostra un nodo solo su una slide specifica
%     es: \node[boxes, on slide=<2>{opacity=0}] {testo};
%
% MACRO:
%   \highlightprofilecolor — colore corrente highlight (ridefinibile)
\newcommand{\highlightprofilecolor}{LightGoldenRod!100}

\tikzset{
	entity/.style={
			draw,
			ellipse,
			very thin
		},
	artifact/.style={
			rectangle,
			draw,
			rounded corners,
			fill=VeryLightGrey!50
		},
	service/.style={
			chamfered rectangle,
			chamfered rectangle angle=30,
			draw,
			%fill=VeryLightGrey!120
		},
	every node/.style={
			font={\scriptsize},
			text centered,
		},
	on link/.style={
			fill=white,
		},
	connector/.style={
			>=stealth,
		},
	connector assurance/.style={
			>=Stealth,
			dashed
		},
	on slide/.code args={<#1>#2}{%
			\only<#1>{\pgfkeysalso{#2}}%
		},
	highlight base/.style={
			very thick
		},
	highlight/.style={
			highlight base,
			fill=LightGoldenRod,
		},
	highlight profile/.style={
			\highlightprofilecolor
		},
	highlight double/.style={
			double=LightGoldenRod
		},
	highlight text/.style={
			text=Goldenrod,
		},
	highlight link/.style={
			%thick,
			color=Goldenrod
		},
	highlight light/.style={
			fill=LightGoldenRod!50
		},
	highlight very light/.style={
			fill=LightGoldenRod!20
		},
	element/.style={
			inner sep=1.5pt,
			font=\footnotesize
		},
	connector/.style={
			>=Stealth,
		},
	lattice layer label/.style={
			text width=2.9cm,
			font={\footnotesize},
		},
	not chosen color/.style={
			black!50
		},
	lattice layer separator/.style={
			dashed,
		},
	service label/.style={
			fill=white,
			circle,
			draw,
			very thick,
			inner sep=2.5pt,
		},
	lattice element bound/.style={
			fill=black!15,
			double
		},
	lattice element selected/.style={
			fill=Goldenrod!15,
			double
		},
	lattice element not chosen/.style={
			not chosen color,
			densely dotted
		},
	lattice element not chosen link/.style={
			not chosen color,
			densely dotted
		},
	service element not chosen/.style={
			text=black!50,
			draw=black!50,
			dashed,
		},
	header/.style={
		},
	main box/.style={
			font={\small}
		},
	link/.style={
			>=Latex,
		},
	activity/.style={
			fill=white
		},
	% for cert history
	arrow/.style={
			draw,
			minimum height=1cm,
			inner sep=2pt,
			shape=signal,
			signal from=west,
			signal to=east,
			signal pointer angle=110,
			top color=DarkBlue,
			bottom color=DarkBlue,
			text=White,
			text width=.24\textwidth,
			font = {\small}
		},
	desc/.style={
			%text width=3.9cm,
			text width=.22\textwidth,
			fill=VeryLightGrey,
			%minimum height=1cm,
			%anchor=south
			% text depth=1cm,
			rectangle,
			font = {\footnotesize}
		},
	link comment/.style={
			dashed
		}
}

\maketitle

\footlinecolor{maincolor}

% ==================///==================///==================///
% ==================/// BODY'S PRESENTATION
% ==================///==================///==================///



%\begin{frame}{Table of Contents}
%	\tableofcontents
%\end{frame}
% !TEX root = ../main.tex
% PATTERN FRAME DISPONIBILI (usati nella triennale):
%
% Frame con colonne testo 70% / icona 30%:
%   \begin{frame}{Titolo}
%     \begin{columns}
%       \begin{column}{.7\textwidth}\begin{itemize}\item \alert{punto}\end{itemize}\end{column}
%       \begin{column}{.2\textwidth}\includegraphics[scale=.2]{assets/icon/nome.png}\end{column}
%     \end{columns}
%   \end{frame}
%
% Frame con enumerate nidificate:
%   \begin{enumerate}\item \alert{Step}\begin{enumerate}\item \textbf{Sub}: desc\end{enumerate}\end{enumerate}
%
% Frame fullscreen image:
%   \includegraphics[width=\textwidth,height=0.9\textheight,keepaspectratio]{assets/file.pdf}
%
% Slide divisore capitolo (da chapters/special_slides.tex):
%   \begin{chapter}[assets/background_negative]{maincolor}{Titolo}\end{chapter}
%   Colori: maincolor, stataledarkgreen, statalered, stataleyellow, statalelilla
%
% Slide con immagine laterale:
%   \begin{sidepic}{assets/path/image}{Titolo}\end{sidepic}

\begin{frame}{}
    % TODO
\end{frame}

\input{sections/02_motivazioni}
% !TEX root = ../main.tex
\begin{frame}{}
    % TODO
\end{frame}

\input{sections/04_metodologia}
\input{sections/05_implementazione}
\input{sections/06_conclusion}



% ==================///==================///==================///
% ==================/// END PRESENTATION
% ==================///==================///==================///

\backmatter

\thispagestyle{empty}
\mbox{}
\clearpage
\setcounter{page}{1}

\appendix
% !TEX root = ../main.tex
% Slide di backup per le domande della commissione.
% Non vengono proiettate durante la presentazione principale.

\begin{frame}{Architettura di APIGuard Assurance}
    \centering\resizebox{\textwidth}{!}{
        % Figura 3-architettura-apiguard adattata (senza float, caption, label, \gls)
        \begin{tikzpicture}[
                node distance=1.5ex and 1.5em,
                base/.style={draw, rounded corners=1.5pt, inner sep=5pt, align=center, font=\small},
                col-in/.style={base, fill=gray!10, draw=gray!50, minimum width=10em},
                col-mid-blue/.style={base, fill=blue!7, draw=blue!40, minimum width=13em, minimum height=3em},
                col-mid-green/.style={base, fill=green!7, draw=green!45, minimum width=13em, minimum height=3em},
                col-mid-gold/.style={base, fill=orange!10, draw=orange!50, minimum width=13em, minimum height=2.5em},
                col-mid-violet/.style={base, fill=violet!7, draw=violet!40, minimum width=13em, minimum height=3em},
                col-out/.style={base, fill=gray!8, draw=gray!50, minimum width=8.5em},
                arr/.style={-latex, semithick, gray!80!black},
                darr/.style={-latex, semithick, dashed, gray!80!black},
                barr/.style={latex-latex, semithick, gray!80!black},
                lbl/.style={font=\footnotesize\itshape, text=gray!70!black, inner sep=3pt},
                bg-engine/.style={draw=blue!30, fill=blue!3, dashed, rounded corners=3pt, inner sep=10pt}
            ]

            \node[col-mid-blue] (tc) {\textbf{\texttt{TargetContext}} {\footnotesize[frozen]}\\{\footnotesize base\_url | attack\_surface | creds}};
            \node[col-mid-green, below=of tc] (tx) {\textbf{\texttt{TestContext}} {\footnotesize[mutable]}\\{\footnotesize token | teardown | shared\_data}};
            \node[col-mid-gold, below=of tx] (dag) {\textbf{DAG Scheduler}\\{\footnotesize ordinamento, batch sequenziali}};

            \begin{scope}[on background layer]
                \node[bg-engine, fit=(tc)(tx)(dag)] (engine) {};
            \end{scope}
            \node[font=\footnotesize\bfseries, text=blue!70!black, anchor=south east, shift={(-1em, 0.5ex)}] at (engine.north) {Assessment Engine};

            \node[col-in, anchor=north east] (cfg) at ([xshift=-1.5em]tc.west |- engine.north) {\textbf{\texttt{config.yaml}}\\{\footnotesize credenziali, URL, param.}};
            \node[col-in, anchor=east] (oas) at ([xshift=-1.5em, yshift=-7ex]tc.west) {\textbf{\texttt{openapi.yaml}}\\{\footnotesize spec formale target}};

            \node[col-out, right=4em of tc] (kong) {\textbf{Kong Gateway}\\{\footnotesize :8000 / :8443 / :8001}};
            \node[col-out, below=8.5ex of kong] (forgejo) {\textbf{Forgejo}\\{\footnotesize target applicativo}};

            \node[col-mid-violet, below=4ex of dag] (ev) {\textbf{EvidenceStore} {\footnotesize[JSONL]}\\{\footnotesize \texttt{evidence.json} | \texttt{report.html}}};

            \draw[arr] (cfg.east) -- (tc.west |- cfg.east);
            \draw[arr] (oas.north) |- ([yshift=-2ex]tc.west);

            \path (tc.east) ++(0, 1.5ex) coordinate (probe-start);
            \path (kong.west) ++(0, 1.5ex) coordinate (probe-end);
            \draw[arr] (probe-start) -- node[pos=0.6, above, align=center, lbl] {probe\\HTTP} (probe-end);

            \draw[barr] (kong.south) -- (forgejo.north);
            \draw[arr] (dag.north) -- (tx.south);
            \draw[arr] (dag.south) -- (ev.north);

            \path (tx.east) ++(0, -1ex) coordinate (admin-start);
            \path (kong.west) ++(0, -1.5ex) coordinate (admin-end);
            \path (admin-start) ++(1.2em, 0) coordinate (admin-turn1);
            \coordinate (admin-turn2) at (admin-turn1 |- admin-end);
            \draw[darr] (admin-start) -- (admin-turn1) -- (admin-turn2) -- (admin-end);
            \node[lbl, align=left, anchor=north west] at ([xshift=0.1em, yshift=-1.5ex]admin-turn2) {Admin API\\{\footnotesize WHITE\_BOX}};
        \end{tikzpicture}
    }% fine \resizebox
\end{frame}

\begin{frame}{Pipeline a sette fasi}
    \centering\resizebox{!}{0.82\textheight}{
        \begin{tikzpicture}[
                node distance=3.2ex and 4em,
                base/.style={draw, rounded corners=1.5pt, inner sep=3pt, align=center, font=\small, minimum height=3.5ex},
                phase-startup/.style={base, fill=blue!7, draw=blue!40, minimum width=13.5em},
                phase-exec/.style={base, fill=green!7, draw=green!45, minimum width=13.5em},
                err-stop/.style={base, fill=red!6, draw=red!50, minimum width=11.5em},
                err-warn/.style={base, fill=orange!8, draw=orange!50, minimum width=11.5em},
                term/.style={draw=gray!50, fill=gray!15, ellipse, minimum height=2.5ex, inner sep=3pt, align=center, font=\footnotesize\ttfamily},
                lbl/.style={font=\footnotesize\itshape, text=gray!70!black, inner sep=2pt},
                arr-ok/.style={-latex, semithick, gray!70!black},
                arr-ko/.style={-latex, semithick, red!60!black},
                arr-warn/.style={-latex, semithick, dashed, orange!80!black},
                bg-startup/.style={draw=blue!30, fill=blue!3, dashed, rounded corners=3pt, inner sep=8pt},
                bg-exec/.style={draw=green!35, fill=green!2, dashed, rounded corners=3pt, inner sep=8pt}
            ]

            \node[term] (start) {apiguard run};
            \node[phase-startup, below=3.5ex of start] (p1) {\textbf{Fase 1:} Inizializzazione};
            \node[phase-startup, below=of p1] (p2) {\textbf{Fase 2:} OpenAPI Discovery};
            \node[phase-startup, below=of p2] (p3) {\textbf{Fase 3:} Costruzione Contesti};
            \node[phase-startup, below=of p3] (p4) {\textbf{Fase 4:} Test Discovery e DAG};
            \node[phase-exec, below=6.5ex of p4] (p5) {\textbf{Fase 5:} Esecuzione (per test)};
            \node[phase-exec, below=of p5] (p6) {\textbf{Fase 6:} Teardown LIFO};
            \node[phase-exec, below=of p6] (p7) {\textbf{Fase 7:} Report};
            \node[term, below=3.5ex of p7] (end) {exit code: 0 / 1 / 2 / 10};

            \node[err-stop, right=4em of p1] (s1) {\texttt{ConfigurationError}\\{\footnotesize avvio interrotto}};
            \node[err-stop, right=4em of p2] (s2) {\texttt{OpenAPILoadError}\\{\footnotesize avvio interrotto}};
            \node[err-stop, right=4em of p4] (s4) {\texttt{DAGCycleError}\\{\footnotesize avvio interrotto}};
            \node[err-warn, right=4em of p5] (w5) {\texttt{TestResult(ERROR)}\\{\footnotesize isolato, pipeline prosegue}};
            \node[err-warn, right=4em of p6] (w6) {\texttt{TeardownError}\\{\footnotesize WARNING strutturato}};

            \begin{scope}[on background layer]
                \node[bg-startup, fit=(p1)(p2)(p3)(p4)] (grp-startup) {};
                \node[bg-exec, fit=(p5)(p6)(p7)] (grp-exec) {};
            \end{scope}

            \node[font=\footnotesize\bfseries, text=blue!70!black, align=right, anchor=east] at ([xshift=-0.5em]grp-startup.west) {Fasi 1-4\\(bloccanti)};
            \node[font=\footnotesize\bfseries, text=green!70!black, align=right, anchor=east] at ([xshift=-0.5em]grp-exec.west) {Fasi 5-7\\(non bloccanti)};

            \draw[arr-ok] (start) -- (p1);
            \draw[arr-ok] (p1) -- node[midway, right, lbl] {OK} (p2);
            \draw[arr-ok] (p2) -- node[midway, right, lbl] {OK} (p3);
            \draw[arr-ok] (p3) -- node[midway, right, lbl] {OK} (p4);
            \draw[arr-ok] (p4) -- node[midway, right, lbl] {OK} (p5);
            \draw[arr-ok] (p5) -- (p6);
            \draw[arr-ok] (p6) -- node[midway, right, lbl] {OK} (p7);
            \draw[arr-ok] (p7) -- (end);

            \draw[arr-ko] (p1.east) -- node[pos=0.55, above, lbl, text=red!70!black] {KO} (s1.west);
            \draw[arr-ko] (p2.east) -- node[pos=0.55, above, lbl, text=red!70!black] {KO} (s2.west);
            \draw[arr-ko] (p4.east) -- node[pos=0.55, above, lbl, text=red!70!black] {KO} (s4.west);
            \draw[arr-warn] (p5.east) -- node[pos=0.55, above, align=center, lbl, text=orange!80!black] {eccezione\\interna} (w5.west);
            \draw[arr-warn] (p6.east) -- node[pos=0.55, above, align=center, lbl, text=orange!80!black] {cleanup\\parziale} (w6.west);
        \end{tikzpicture}
    }% fine \resizebox
\end{frame}

\begin{frame}{Topologia DAG}
    \centering\resizebox{\textwidth}{!}{
        % Figura 4-dag-topology adattata
        \begin{tikzpicture}[
                node distance=1.5ex and 0.8em,
                tst/.style={draw=orange!55, fill=yellow!18, rounded corners=1.5pt, minimum width=4.5em, minimum height=3.5ex, font=\footnotesize\ttfamily, align=center},
                tst-hl/.style={tst, draw=orange!70, fill=yellow!35, thick},
                ext/.style={tst, minimum width=6.5em, font=\scriptsize\ttfamily},
                phB/.style={draw=blue!40, fill=blue!7, rounded corners=1.5pt, minimum width=5.5em, minimum height=3.5ex, font=\footnotesize\ttfamily, align=center},
                bg-A/.style={draw=orange!50, fill=yellow!6, dashed, rounded corners=3pt, inner sep=12pt},
                bg-B/.style={draw=blue!40, fill=blue!4, dashed, rounded corners=3pt, inner sep=10pt},
                arr/.style={-latex, semithick, blue!60},
                darr/.style={-latex, semithick, dashed, gray!60}
            ]

            \node[tst] (d0a) {0.1};
            \node[tst, right=of d0a] (d0b) {0.2};
            \node[tst, right=of d0b] (d0c) {0.3};
            \node[tst-hl, below=of d0a] (a11) {\textbf{1.1}};
            \node[tst, right=of a11] (d1b) {1.5};
            \node[tst, right=of d1b] (d1c) {1.6};
            \node[tst, below=3.5ex of a11] (d3a) {3.3};
            \node[tst, below=of d3a] (d4a) {4.1};
            \node[tst, right=of d4a] (d4b) {4.2};
            \node[tst, right=of d4b] (d4c) {4.3};
            \node[tst, below=of d4a] (d6a) {6.2};
            \node[tst, right=of d6a] (d6b) {6.4};
            \node[tst, below=of d6a] (d7a) {7.2};
            \node[ext, below=of d7a.south west, anchor=north west] (ex0) {ext.0.1.nuclei};
            \node[ext, below=of ex0.south west, anchor=north west] (ex1) {ext.1.5.sslyze};
            \node[ext, right=of ex1] (ex2) {ext.1.5.testssl};

            \begin{scope}[on background layer]
                \node[bg-A, fit=(d0a)(d0c)(a11)(d1c)(d4c)(d7a)(ex1)(ex2)] (phaseA) {};
            \end{scope}
            \node[font=\footnotesize\bfseries, text=orange!80!black, anchor=south west] at ([xshift=1ex, yshift=0.5ex]phaseA.north west) {Phase A (16 test, nessuna dipendenza)};

            \path (a11.south) ++(0, -1.75ex) coordinate (bus-y);
            \path (phaseA.east |- bus-y) coordinate (exitA);
            \path (exitA) ++(2.5em, 0) coordinate (split);

            \node[phB, anchor=west] (b14) at ([xshift=1.5em, yshift=3ex]split) {1.4};
            \node[phB, anchor=west] (b21) at ([xshift=1.5em, yshift=-3ex]split) {2.1};

            \begin{scope}[on background layer]
                \node[bg-B, fit=(b14)(b21)] (phaseB) {};
            \end{scope}
            \node[font=\footnotesize\bfseries, text=blue!80!black, anchor=south west] at ([xshift=1ex, yshift=0.5ex]phaseB.north west) {Phase B (\texttt{depends\_on=["1.1"]})};

            \draw[blue!60, semithick] (a11.south) |- (exitA) -- node[above, yshift=0.3ex, font=\scriptsize\itshape, text=gray!80!black, fill=white, inner sep=0.1pt] {\texttt{1.1} $\rightarrow$} (split);
            \draw[arr] (split) |- (b14.west);
            \draw[arr] (split) |- (b21.west);
        \end{tikzpicture}
    }% fine \resizebox
\end{frame}

\begin{frame}{Mappa di copertura (v0.1.0)}
    \centering\resizebox{!}{0.82\textheight}{
        \begin{tikzpicture}[
                node distance=1.5ex and 0.4em,
                base/.style={draw, rounded corners=1.5pt, align=center, inner sep=4pt},
                dom/.style={base, fill=gray!10, draw=gray!50, minimum width=10em, minimum height=5.2ex, font=\footnotesize\bfseries},
                pass/.style={fill=green!10, draw=green!50!black},
                fail/.style={fill=red!10, draw=red!60!black},
                skip/.style={fill=gray!15, draw=gray!50},
                tst/.style={base, minimum width=3.5em, minimum height=5.2ex, font=\footnotesize\ttfamily},
                ext/.style={base, minimum width=5em, minimum height=5.2ex, font=\footnotesize\ttfamily\itshape},
                inact/.style={base, draw=gray!35, dashed, fill=gray!4, minimum width=8em, minimum height=5.2ex, font=\footnotesize\itshape, text=gray!70!black},
            ]

            \node[dom] (d0) {D0: Discovery};
            \node[tst, fail, right=1.5em of d0] (t01) {\texttt{0.1}};
            \node[tst, fail, right=of t01] (t02) {\texttt{0.2}};
            \node[tst, skip, right=of t02] (t03) {\texttt{0.3}};
            \node[ext, pass, right=of t03] (te01) {\texttt{ext.0.1}\\[-2pt]{\scriptsize nuclei}};

            \node[dom, below=1.5ex of d0.south west, anchor=north west] (d1) {D1: Authentication};
            \node[tst, fail, right=1.5em of d1] (t11) {\texttt{1.1}};
            \node[tst, pass, right=of t11] (t14) {\texttt{1.4}};
            \node[tst, pass, right=of t14] (t15) {\texttt{1.5}};
            \node[tst, skip, right=of t15] (t16) {\texttt{1.6}};
            \node[ext, fail, right=of t16] (te15s) {\texttt{ext.1.5}\\[-2pt]{\scriptsize sslyze}};
            \node[ext, fail, right=of te15s] (te15t) {\texttt{ext.1.5}\\[-2pt]{\scriptsize testssl}};

            \node[dom, below=1.5ex of d1.south west, anchor=north west] (d2) {D2: Authorization};
            \node[tst, pass, right=1.5em of d2] (t21) {\texttt{2.1}};

            \node[dom, below=1.5ex of d2.south west, anchor=north west] (d3) {D3: Data Integrity};
            \node[tst, pass, right=1.5em of d3] (t33) {\texttt{3.3}};

            \node[dom, below=1.5ex of d3.south west, anchor=north west] (d4) {D4: Availability};
            \node[tst, fail, right=1.5em of d4] (t41) {\texttt{4.1}};
            \node[tst, pass, right=of t41] (t42) {\texttt{4.2}};
            \node[tst, pass, right=of t42] (t43) {\texttt{4.3}};

            \node[dom, fill=gray!4, draw=gray!35, below=1.5ex of d4.south west, anchor=north west] (d5) {D5: Observability};
            \node[inact, right=1.5em of d5] (td5) {nessun test attivo};

            \node[dom, below=1.5ex of d5.south west, anchor=north west] (d6) {D6: Hardening};
            \node[tst, pass, right=1.5em of d6] (t62) {\texttt{6.2}};
            \node[tst, pass, right=of t62] (t64) {\texttt{6.4}};

            \node[dom, below=1.5ex of d6.south west, anchor=north west] (d7) {D7: Business Logic};
            \node[tst, fail, right=1.5em of d7] (t72) {\texttt{7.2}};

            \node[tst, pass, below=3.5ex of d7.south west, anchor=north west] (lg-pass) {\texttt{PASS}};
            \node[tst, fail, right=0.8em of lg-pass] (lg-fail) {\texttt{FAIL}};
            \node[tst, skip, right=0.8em of lg-fail] (lg-skip) {\texttt{SKIP}};
            \node[ext, pass, right=0.8em of lg-skip] (lg-ext) {\textit{Esterno}};
        \end{tikzpicture}
    }% fine \resizebox
\end{frame}

\end{document}


